% !PRINTED 431
\weberpage{431}

\S\ 119. \gesperrt{The Discriminant of the Invariant Equation.}

We now raise the question when an invariant equation
\weq{1}{F_p(v,\,u) = 0,}
in which $p$ is a prime, has two or more equal roots $v$ for $u = j(\omega)$, or when among the $p+1$ quantities
\weq{2}{j(p\,\omega),\quad j\!\left(\frac{\omega+c}{p}\right),\quad c = 0,\,1,\,\ldots,\,p-1}
two or more occur that are equal to one another. This takes place then and only then, when among the $p+1$ quantities
\weq{3}{p\,\omega,\quad \frac{\omega+c}{p}}
two equivalent ones occur. Accordingly $\omega$ must in any case satisfy a quadratic equation with integral coefficients and negative discriminant $D=-m$, which we assume in the form:
\weq{4}{A\,\omega^2 + B\,\omega + C = 0,\quad 4\,A\,C - B^2 = m > 0,}
in which $A,\,B,\,C$ are without a common divisor.

If we replace $\omega$ by an equivalent number, and hence (4) by an equivalent equation, the quantities (2) are only interchanged among themselves. The question as to the \gesperrt{number} of equal roots of (1) is therefore not affected by this, and we may accordingly assume that $A$ is not divisible by $p$. If now, in the first place, $p\,\omega$ is equivalent to $\dfrac{\omega+c}{p}$, then:
\weq{5}{\frac{\omega+c}{p} = \frac{\gamma+\delta\,p\,\omega}{\alpha+\beta\,p\,\omega},\quad \alpha\,\delta-\beta\,\gamma = 1.}
If we write this equation thus:
\weq{6}{\beta\,p\,\omega^2 + (\alpha+\beta\,c\,p-\delta\,p^2)\,\omega + c\,\alpha-\gamma\,p = 0,}
then it follows by comparison with (4) (since $A$ was to be indivisible by $p$) that $\alpha$ is divisible by $p$, and must therefore be $=p\,\alpha'$, and that an integer $x$ exists which satisfies the conditions:
\weq{7}{\beta = A\,x,\quad \alpha'+\beta\,c-\delta\,p = B\,x,\quad c\,\alpha'-\gamma = C\,x.}
If we further put
\weq{8}{\alpha'-\beta\,c+\delta\,p = y,}
then it follows:
\weq{9}{\begin{aligned}
2\,\alpha' &= B\,x+y,\\
2\,(\beta\,c-\delta\,p) &= B\,x-y.
\end{aligned}}
%CONT>>
